\subsection{Between Inequalities} 
When we have certain compound inequalities which use \makeblank[0.5in], we can string them together to say that one number is ``between'' two others.
\begin{Example}
We can write
\[
4<2x-8\leq 12
\]
to mean
\[
\makeblank[1.5in]
\mbox{ and }
\makeblank[1.5in]
\]
\end{Example}
\begin{Note}
We \emph{cannot} use this notation for a compound inequality with \makeblank[0.5in], ever. When we write this ``between'' type inequality, it \emph{always} means that \emph{both} inequalities must be true.
\end{Note}
We can solve a ``between'' inequality in two ways.
\begin{Example}
Solve $4<2x-8\leq 12$
\begin{lpic}{}{7}
\twocol{-1}{7}{Method 1:}{Method 2:}
\end{lpic}
\end{Example}
We can only use the the method of working with ``all three sides'' if the variable only appears in the middle.
\begin{Example}
Solve $x+3<9<4x-7$
\begin{cpic}{}{6}
\end{cpic}
\end{Example}
Note that the ``between'' inequality represents a graph in the number line with \makeblank[1in] (empty or filled) at both ends.